\documentclass[12pt]{article}
\usepackage[margin=1in]{geometry}
\usepackage{setspace}
\usepackage{parskip}
\setstretch{1.2}

\title{Informe general de la Plataforma Fellowship}
\author{Equipo de desarrollo}
\date{\today}

\begin{document}
\maketitle
\tableofcontents
\newpage

\section{Resumen}
La Plataforma Fellowship es un espacio digital para organizar proyectos de estudiantes. Su idea principal es sencilla: reunir en un solo lugar el registro de equipos, la aprobación de nuevos integrantes, la entrega de tareas y la solicitud de reservas. De este modo, las personas que administran el programa pueden ver todo el avance sin usar hojas de cálculo ni mensajes dispersos.

Este informe explica su funcionamiento con palabras simples. No se enfoca en tecnología, sino en las personas y en el flujo de trabajo. La idea es que cualquier lector entienda qué hace la plataforma, cómo se usa y por qué aporta orden y claridad.

\section{Qué es y para qué sirve}
La plataforma funciona como una oficina virtual para equipos y administradores. Los equipos pueden registrar su proyecto, seguir sus tareas y comunicarse con el programa. Los administradores pueden revisar solicitudes, aprobar equipos, dar retroalimentación y coordinar el uso de espacios como laboratorios.

En palabras simples: ayuda a ordenar el trabajo, evita confusiones y permite que todos sepan en qué etapa está cada proyecto.

También cumple una función de memoria. Si una persona nueva se integra al equipo, puede ver el historial y ponerse al día. Si un administrador cambia, el programa no se detiene porque la información ya está organizada.

\section{Quiénes la usan}
\begin{itemize}
  \item \textbf{Estudiantes y equipos}: registran su proyecto, suben entregas y solicitan apoyo.
  \item \textbf{Administradores}: revisan solicitudes, validan equipos, asignan tareas y llevan el control general.
\end{itemize}

En algunos casos también participan mentores o asesores. Aunque no tengan acceso directo a la plataforma, su trabajo se refleja en comentarios, retroalimentación y decisiones que quedan registradas.

\section{Cómo funciona, paso a paso}
\textbf{Paso 1: Registro}. Un equipo crea su proyecto o se une con un código de invitación. La solicitud queda en espera.

\textbf{Paso 2: Revisión}. Un administrador revisa la solicitud y decide si la aprueba.

\textbf{Paso 3: Trabajo diario}. Con la aprobación, el equipo accede a su tablero, ve las tareas pendientes y entrega evidencias. Si necesita una asesoría o un laboratorio, puede solicitar una reserva.

\textbf{Paso 4: Seguimiento}. Los administradores revisan las entregas, dejan comentarios y ven el avance general del programa.

Este flujo busca evitar pasos innecesarios. Primero se valida quién entra, luego se facilita el trabajo cotidiano. El equipo no necesita preguntar todo el tiempo qué sigue o dónde dejar una tarea: la plataforma ya lo muestra.

\section{Ejemplo sencillo de uso}
Imaginemos un equipo llamado ``Brillo Solar''. Sus integrantes se registran y esperan aprobación. Un administrador revisa su información y les da acceso. Desde ese momento:
\begin{itemize}
  \item Ven una lista clara de tareas con fechas y prioridades.
  \item Suben un enlace con su avance semanal.
  \item Reciben un comentario cuando la entrega fue revisada.
  \item Solicitan una sesión en el laboratorio y esperan confirmación.
\end{itemize}

En pocas palabras, el equipo sabe qué hacer y el administrador puede dar seguimiento sin buscar mensajes en varios canales.

\section{Beneficios principales}
\begin{itemize}
  \item \textbf{Orden}: toda la información del programa queda en un solo lugar.
  \item \textbf{Claridad}: cada equipo sabe qué debe hacer y en qué etapa está.
  \item \textbf{Control}: el proceso de aprobación evita registros no autorizados.
  \item \textbf{Seguimiento}: administradores y equipos pueden ver avances y pendientes.
  \item \textbf{Coordinación}: es más sencillo pedir y aprobar reservas de espacios.
\end{itemize}

Además, hay beneficios menos visibles pero igual de importantes: se reduce el tiempo perdido en coordinaciones, se evitan confusiones sobre entregas y se mejora la comunicación entre equipos y administradores.

\section{Lo que no es la plataforma}
La plataforma no reemplaza el trabajo creativo ni las reuniones entre personas. Tampoco es una red social. Su propósito es más práctico: organizar, registrar y hacer seguimiento. El valor está en el orden y en la claridad, no en la complejidad.

\section{Responsabilidades de cada rol}
\textbf{Equipos}:
\begin{itemize}
  \item Mantener su información actualizada.
  \item Entregar evidencias con claridad.
  \item Respetar fechas y solicitudes de reserva.
\end{itemize}

\textbf{Administradores}:
\begin{itemize}
  \item Revisar solicitudes y aprobar con criterio.
  \item Dar retroalimentación oportuna.
  \item Mantener el calendario de reservas organizado.
\end{itemize}

\section{Buenas prácticas de uso}
\begin{itemize}
  \item Revisar el tablero al inicio de cada semana.
  \item Subir entregas completas y con un título claro.
  \item Leer los comentarios antes de enviar una nueva versión.
  \item Solicitar reservas con tiempo suficiente.
  \item Evitar duplicar solicitudes o tareas.
\end{itemize}

Estas prácticas simples reducen retrasos y hacen que la plataforma sea realmente útil.

\section{Seguridad y cuidado de la información}
La plataforma tiene un proceso de aprobación antes de que un equipo pueda trabajar. Esto ayuda a mantener el orden y a evitar accesos no deseados. También se guarda un historial básico de acciones para que el programa pueda revisar lo ocurrido si hay dudas.

Es importante que los usuarios cuiden sus credenciales y que no compartan su acceso con terceros. La plataforma organiza la información, pero el buen uso depende de cada persona.

\section{Preguntas frecuentes}
\textbf{¿Qué pasa si un equipo se equivoca al registrarse?} Se puede corregir la información desde el panel, con apoyo del administrador.

\textbf{¿Qué sucede si una entrega se sube tarde?} Queda registrada con su fecha real, lo que permite a los administradores tomar decisiones claras.

\textbf{¿Se pueden crear equipos nuevos en cualquier momento?} Sí, siempre que la administración lo autorice.

\textbf{¿Qué ocurre si una reserva no se aprueba?} El equipo puede solicitar otra fecha o ajustar su plan.

\section{Conclusiones}
La Plataforma Fellowship es una herramienta pensada para hacer más simple la gestión de proyectos estudiantiles. Centraliza tareas, aprobaciones y reservas, y permite que el programa funcione con menos fricciones. En resumen: ayuda a que los equipos trabajen mejor y a que los administradores tengan una visión clara del avance.

\section{Anexo: glosario simple}
\begin{itemize}
  \item \textbf{Proyecto}: el trabajo que desarrolla un equipo.
  \item \textbf{Equipo}: grupo de estudiantes que trabaja en un proyecto.
  \item \textbf{Entrega}: evidencia o avance que el equipo presenta.
  \item \textbf{Aprobación}: decisión para permitir el ingreso de un equipo.
  \item \textbf{Reserva}: solicitud de uso de un espacio o asesoría.
\end{itemize}

\end{document}
